% !TEX TS-program = lualatex
\documentclass[../ctenarsky-denik.tex]{subfiles}
\begin{document}
\kniha{
  nazev=Hra o trůny,
  serie=Píseň ledu a ohně,
  dil=1,
  podtitul=První kniha ságy,
  autor=George R. R. Martin,
  obalka=hra-o-truny,
  hodnoceni=5,
  precteno=listopad 2026,
}


\podnadpis{Děj}
První román začíná po asi patnácti letech Robertovy vlády. Prolog se odehrává na sever od obří Zdi, která má za úkol chránit Sedm království před nebezpečím ze severu. Skupina mužů z Noční hlídky, kteří mají za úkol Zeď chránit, je na stopě lidí, kterým říkají Divocí. Jsou však zmasakrováni nemrtvými tvory, Jinými. Jediného přeživšího ze skupiny pak jako dezertéra z Noční hlídky popraví lord Eddard Stark, strážce jednoho z království – severu.

Lorda Starka po dlouhé době navštíví jeho přítel, král Robert Baratheon, a žádá ho, aby se ujal úřadu králova pobočníka v hlavním městě říše, Králově přístavišti, poté, co zemřel jeho předchozí pobočník Jon Arryn. Stark má své pochybnosti, ale nakonec souhlasí, zčásti i proto, že z dopisu od Arrynovy manželky (sestry jeho manželky Catelyn) ví, že Arryn byl zavražděn, a chce zjistit, jaké okolnosti k tomu vedly.

Lord Stark do Králova přístaviště bere i své dvě dcery, z nichž starší Sansa je snoubenkou prince Joffreyho. Nakonec zjistí i to, co stálo život Jona Arryna. Otcem tří dětí královny Cersei Lannister není král Robert, ale narodily se z incestu s jejím bratrem-dvojčetem Jaimem. Povaha Eddarda Starka mu velí Cersei varovat a dát jí možnost s dětmi utéct před Robertovým hněvem, než mu řekne pravdu. Král Robert však vzápětí umírá na zranění z lovu a zatímco Eddard Stark se snaží zasadit se o právo Robertova skutečného dědice (jeho bratra Stannise Baratheona), Cersei ho nechá uvěznit. Eddarda Starka nakonec nechá před očima jeho dcery Sansy popravit nový král Joffrey – proti vůli Cersei i královské rady, než mají šanci zakročit.

Mezitím se z Králova přístaviště vydává v přestrojení za chlapce zpět na sever mladší z dcer Arya Stark. Dědic Eddarda Starka, patnáctiletý Robb, je svými vazaly prohlášen králem severu a má v úmyslu osamostatnit své mladé království z moci Králova přístaviště.

Druhá dějová linie ságy se odehrává na Zdi a na sever od ní. Jejím protagonistou je nemanželský syn Eddarda Starka Jon Sníh. V první knize se Jon rozloučí s rodinou, seznamuje se se svými novými bratry z Noční hlídky a na konci se vydává v jejich doprovodu za Zeď.

Poslední dějová linie popisuje osud dědičky rodu Targaryenů, která musela ještě jako dítě prchnout do exilu a nyní se postupně připravuje na svůj návrat do Západozemí, aby uplatnila svůj nárok na trůn. Její psychicky labilní bratr Viserys ji prodá Khalu Drogovi, vůdci divokých Dothraků, za příslib armády. Viseryse ovšem jeho netrpělivost stojí život, když ho Khal Drogo zabije. Daenerys nakonec přijde i o svého manžela Droga a zůstává jen s hrstkou věrných a čerstvě narozenými draky.


\podnadpis{Ukázka}
\begin{ukazka}
  „Milosrdenství nikdy není chybou, lorde Renly,“ opáčil Ned. „Na Trojzubci tehdy ser Barristan skolil tucet dobrých mužů, Robertových přátel a mých. Když ho před nás přivedli, vážně zraněného a blízko smrti, Roose Bolton na nás naléhal, aby mu nechali podříznout hrdlo, ale tvůj bratr řekl: ‚Nezabiju muže za věrnost, za to, že dobře bojoval,‘ a poslal pro svého vlastního mistra, aby seru Barristanovi ošetřil rány.“ Pohlédl na krále dlouhým chladným pohledem. „Jinak by tady s námi dnes tento muž neseděl.“
  \zdrojukazky{kapitola I}
\end{ukazka}

\muj{Kniha se mi líbila, protože\ldots}
\end{document}
